% !TeX root = ../papers/nc-p3.tex
\section{四元群窓の閉軌道への全退化とトーリック商}
\label{sec:klein-polystable-limits}
$W$ を総文字零の64語の線形空間、$T$ を四位置の対角トーラスの射影像、
$G=\prod\operatorname{PGL}(V_i)$ とする。射影GITには
$\prod\operatorname{SL}(V_i)$ の自然な線形化を用いる。以下の閉軌道は半安定集合内で取る。
\begin{theorem}\label{thm:klein-polystable-limits}
正則表現型四元群を持つ全標準AS窓は半安定であり、その軌道閉包にある唯一の閉軌道も
標準Koszul左右AS正則である。従って同じ印を持つ二隣接階数6の候補が正則であることと、
半安定で閉軌道の窓が正則であることは同値である。閉軌道は次の有限操作で求まる。
\begin{enumerate}
\item 両切断に欠色があれば、四文字が相異なる語は全て残す。
反復対 $ijij,jiji$ は、その補対 $klkl,lklk$ がない場合だけ二つとも消す。
\item 疎な平行二重色の正則台で、番号
$51,54,55,56,57,58,62,63,64,65,76,78,112,115,117$ は指定始点を保つ
文字置換を含めて台50に移り、$59,66$ は台52に移る。
残す語の係数はそのまま使う。
\item 他の全正則表示は既に閉軌道である。
\end{enumerate}
さらに四元群印を仮定しない二隣接階数6の任意の窓も、ここで得た正則な閉軌道のいずれかを
極限に持てば標準Koszul左右AS正則である。
\end{theorem}
\begin{proof}[台の多面体と四位置での閉軌道]
語 $q$ の重みを前節の $E$ の行とし、各位置の四座標が全て $1/4$ の点を $o$ と書く。
台 $S$ の重みの凸包を $P$ とする。$o\notin P$ なら対角一径数群が零へ送る。
$o\in P$ なら $o$ を相対内部に含む最小面 $F$ があり、外の語を消す
四つの跡零整数重みを一つ取れる。残った台には、全語で正かつ各位置の周辺和が全て等しい
有理重みがある。このとき対角トーラスでの軌道は多安定である。
複素数上では実対角倍率によるノルム最小化で全ての一位置周辺ノルムが等しくなる。
実際、対数変数での二乗ノルムは重みの指数関数の正係数和で、台が張る方向では
$o$ が相対内部にあるので無限遠で発散する。
各位置の異なる二文字の切片の台は総文字零により交わらず、Gram行列の非対角成分は元から零である。
従って全 $G$ に対しても最小ノルム条件を満たし、閉軌道を得る。
有限個の係数の体を複素数へ埋め込んでから代数閉拡大に戻せる。

$F$ が真の面なら元の窓はこの閉軌道と同型にならない。実際、一つの $G$ 軌道内で
$K$ が固定する射影点の部分の接空間は接表現の不変部分であり、
$\operatorname{Lie}G$ から軌道接空間への全射を $K$ 不変に取っても全射である。
よって各固定点で中心化群の恒等成分の軌道が開である（固定部分は滑らかで、
または接次元と軌道次元が等しいことを使える）。その恒等成分は $T$ なので各軌道は
固定部分で閉でもある。相異なる台の対角極限が元の $G$ 軌道に入ることはない。
\end{proof}

\begin{proof}[正則な台の全場合]
両欠色では節\ref{sec:unmarked-klein-pbw}の二反転により、全相異なる語は同じ $(b,e)$ の
四項が同時に現れる。この四項は各位置で全文字を一度ずつ使う。
補い合う二反復対の四項も同じ性質を持つ。残る反復対の片割れを消すには、
残った辺の二ラベルに四位置で重み $+1$、その補辺に $-1$ を与える。
全相異なる語、他色の反復対の重みは零であり、残った辺の二項だけが正である。
各欠損色でこれを加える。全相異なる語が全くない場合は零行の不在から全反復対がある。
従って残る台は空でなく、各残存語を正の重みで覆えるのでちょうど最小面である。
一回の片割れ消去は節\ref{sec:unmarked-klein-pbw}の良還元の証明にある階数2の
主元と非零行・列を保つ。その段階で新しく全部欠けた反復色があれば同節のPBW判定が
適用でき、残る片割れを順に消しても同じ主元の議論が使える。

一方だけに欠色の正則台、異方向の疎な二正則台では各台に正の重みを直接与えられる。
$3111$ 型の全非定数60語は各位置の辺際が等しく、その重み行列の階数13は全64語の階数に等しい。
従って定数四語の任意の追加に対し $o$ は相対内部にとどまる。
異方向の完備な薄色48語でも同じ階数13と周辺和の一致が成り立つ。
平行方向の48語は同じ周辺和を持ち、階数12であり、そのアフィン包は高位文字数2の超平面である。
可逆交差部と混合表示には全低位・全高位の語が双方存在する。
二重交換表示でも節\ref{sec:klein-double-exchange-fourier}の第一型には $X\ne Y,W\ne V$、
第二型には $k/g_i,k/r_i\ne0$ があり、両側に語がある。
よって超平面の両側に重みが追加され $o$ は再び相対内部となる。

平行疎型の全正則代表については
\path{computations/verify_klein_polystable_limits.py} が各表示の台を検査する。
定理の17台には消す語で正、残す語で零の整数重みを与え、残存台が実際に50,52の
始点固定の付け替えに属することを確かめる。残りの正則台、50,52自身、上記の片欠色と
異方向の疎台の全てで、正有理ベクトル $c$ により $E_S^{\mathsf t}c=\boldsymbol1$ を検算する。
二つの終着台の各厚色・各切断には非零単項式の二次小行列式があり、全薄色は非零である。
階数上界は特殊化で保存されるので、係数の特殊値を含め定理\ref{thm:klein-parallel-flags}で正則である。
最後に極限が正則なら、元の軌道からその極限を通る曲線を取って
命題\ref{prop:unrestricted-good-reduction}を適用できる。これは元の点の対称性を使わない。
\end{proof}

\begin{theorem}\label{thm:klein-toric-normalization}
$\mathbb P(W)//T$ は次元51の正規射影トーリック多様体で、その前節の384元群による商は
$\mathbb P(W)$ の全基底変更GIT商における像の正規化である。
特に任意の印付き族の多安定な特殊化の像はこの有限写像の下でその閉包を取って得られる。
\end{theorem}
\begin{proof}
四位置の跡零対角群を $\widetilde T$ とする。
$B=k[W]^{\widetilde T}$ は、各位置の四ラベルで周辺次数が等しい非負整数指数の
単項式環である。直交格子と正錐の交わりは飽和しており正規である。
全重みの階数13、$o$ は内部なので $\operatorname{Proj} B$ は次元 $64-13=51$ を持つ。
全 $\prod\operatorname{SL}(V_i)$ の斉次不変式の制限の像を $A\subset B$ とする。
上の最小面の議論は階数条件を使わず、対角群と全基底群の零錐が $W$ 内で一致することを示す。
従って $A_+$ の $B$ での零点は斉次頂点だけである
（アフィン対角商 $W\to\operatorname{Spec}B$ の全射性を使う）。
$B/A_+B$ は有限次元であり、斉次次数の帰納により $B$ は $A$ 上有限である。
これで射影商の像への有限全射が従う。射影群による付け替えは全基底変更商上で自明なので
384元群の商を経由する。

系\ref{cor:klein-four-configurations}の標本は今示した多安定な点で、射影安定群が正確に $K$ である。
有限安定群なので安定点であり、Lunaのエタール切片によりその飽和近傍で各安定群は
この有限群の部分群に共役である。従って $\mathbb P(W)$ の一般点では指定 $K$ が
既に全安定群である。安定点の二つの表示の同値は指定 $K$ を正規化する四位置変換だけから
生じ、前節より対角群と384置換で尽くされる。有限写像はこの商の後で一般に一対一だから
標数零で双有理である。正規環の有限群不変環の正規性で結論が従う。
有限射は閉写像なので、任意の部分集合の像の閉包についての主張も従う。
\end{proof}
